% Figure: oblique-shock geometry on a wedge. Upstream flow M1 turns by
% deflection angle theta across a shock inclined at angle beta.
\begin{figure}[htbp]
\centering
\begin{tikzpicture}[scale=1.0]
% wedge: lower surface horizontal, upper surface inclined at theta=15 deg
% apex at origin, opening to the right
\def\theta{15}
\def\beta{32}
% wedge body (filled gray)
\fill[enggray!20] (0,0) -- (5,0) -- (5,{5*tan(\theta)}) -- cycle;
\draw[engblack, thick] (0,0) -- (5,0);
\draw[engblack, thick] (0,0) -- (5,{5*tan(\theta)});
% incoming flow arrows (horizontal) from the left toward the wedge tip
\foreach \y in {0.5,1.1,1.7} {
  \draw[engblue, -{Latex[length=2mm]}] (-2.3,\y) -- (-0.9,\y);
}
\node[engblue, font=\scriptsize, anchor=east] at (-2.35,1.1) {$M_1$};
% oblique shock line from apex at angle beta above the lower surface
\draw[engred, very thick] (0,0) -- ({4.2*cos(\beta)},{4.2*sin(\beta)});
\node[engred, font=\scriptsize, anchor=south west] at ({4.2*cos(\beta)},{4.2*sin(\beta)}) {shock};
% deflected flow downstream (parallel to upper wedge surface), above wedge
\foreach \y in {2.6,3.1} {
  \draw[engorange, -{Latex[length=2mm]}] ({0.9*\y/3},{\y}) -- ({0.9*\y/3 + 1.3*cos(\theta)},{\y + 1.3*sin(\theta)});
}
\node[engorange, font=\scriptsize, anchor=south] at (2.9,3.4) {$M_2$ (deflected by $\theta$)};
% angle beta arc at apex (between lower surface and shock)
\draw[engred] (0.9,0) arc (0:\beta:0.9);
\node[engred, font=\scriptsize] at ({1.15*cos(\beta/2)},{1.15*sin(\beta/2)}) {$\beta$};
% angle theta arc (between lower surface direction and upper wedge surface)
\draw[engblack] (2.0,0) arc (0:\theta:2.0);
\node[engblack, font=\scriptsize] at (2.4,{2.4*tan(\theta/2)}) {$\theta$};
\end{tikzpicture}
\caption[Oblique-shock geometry on a wedge]{An attached oblique shock
on a wedge. The supersonic upstream flow $M_1$ meets the compression
corner and turns through the deflection angle $\theta$ (the wedge
half-angle) across a shock inclined at the shock angle $\beta$ to the
upstream direction. The velocity component normal to the shock
behaves exactly like a normal shock and is compressed; the component
tangential to the shock is unchanged. For each $M_1$ and each
$\theta$ below the maximum $\theta_{\max}(M_1)$ the
$\theta$=$\beta$=$M$ relation gives two shock angles, a weak shock
(the smaller $\beta$ drawn here) and a strong shock; physical flows
over an unconfined wedge take the weak branch. For $\theta>
\theta_{\max}$ no attached oblique shock exists and the shock detaches
into a curved bow shock standing ahead of the body.}
\label{fig:vol03ch13:oblique-shock}
\end{figure}
